\documentclass[11pt]{article}
\usepackage[a4paper,margin=1in]{geometry}
\usepackage{amsmath,amssymb,amsthm}
\usepackage{graphicx}
\usepackage{xcolor}
\usepackage[hidelinks]{hyperref}
\usepackage{microtype}
\setlength{\emergencystretch}{2em}

\newtheorem{theorem}{Theorem}
\newtheorem{lemma}[theorem]{Lemma}
\newtheorem{corollary}[theorem]{Corollary}
\newcommand{\D}{\mathbb D}
\newcommand{\Htwo}{H^2}
\newcommand{\K}{\mathcal K}
\newcommand{\ran}{\operatorname{ran}}
\newcommand{\clos}{\operatorname{clos}}
\newcommand{\rhoh}{\rho}

\title{Rigidity toward the polynomially bounded shift-similarity problem}
\author{Automated mathematical discovery run \texttt{fa\_banach\_001}}
\date{13 August 2026}

\begin{document}
\maketitle

\begin{center}
\fcolorbox{black}{gray!8}{\parbox{0.91\linewidth}{
\textbf{Status.} Candidate substantial partial result, likely valid, subject
to human review.  The full polynomially bounded problem is not resolved.
We prove two affirmative structural cases, a local Riesz-rigidity theorem,
and a cluster--orthogonality obstruction which rules out the most direct use
of the 2024 cyclic polynomially bounded construction as a counterexample.
}}
\end{center}

\section{Source problem}

Let $S$ be the unilateral shift on $\Htwo$ and
\[
 h_\lambda(z)=\frac1{1-\overline\lambda z},\qquad
 u_\lambda=(1-|\lambda|^2)^{1/2}h_\lambda
 \quad(\lambda\in\D).
\]
Thus $\|u_\lambda\|=1$ and $S^*u_\lambda=\overline\lambda
u_\lambda$.  The source is M.~F.~Gamal', \emph{On power bounded operators
with holomorphic eigenvectors, II}, arXiv:1901.03883~\cite{Gamal2019}.
It asks whether the following contraction criterion of
Uchiyama~\cite{Uchiyama} remains true for polynomially bounded operators.

\medskip
\noindent\textbf{Question.}
Suppose $T\in\mathcal L(\mathcal H)$ is polynomially bounded and there is a
quasiaffinity $X:\mathcal H\to\Htwo$ such that $XT=SX$ and
\begin{equation}\label{eq:kernel-estimate}
 0<c\leq \|X^*u_\lambda\|\leq C<\infty
 \qquad(\lambda\in\D).
\end{equation}
Must $T$ be similar to $S$?

The source says explicitly that the question remains open; the relevant
passage on source p.~2 is reproduced below.

\begin{center}
\includegraphics[width=\linewidth]{figures/open_problem_crop.png}
\end{center}

The 2019 paper disproves the power-bounded version even after requiring
$T\sim S$.  A later paper~\cite{Gamal2024} constructs cyclic polynomially
bounded operators not similar to contractions and quasisimilar to $S$.
The latter examples are an immediate candidate for settling the displayed
question negatively, but they do not state the estimate
\eqref{eq:kernel-estimate}.  Section~\ref{sec:audit} proves that their direct
shift extension cannot satisfy it.

\section{Proof intuition}

There are two competing geometries.  Polynomial boundedness lets bounded
analytic interpolation act diagonally on the eigenvectors $X^*u_\lambda$.
On every interpolating sequence this forces those normalized eigenvectors to
be a Riesz sequence.  This is genuine rigidity, and it proves the desired
conclusion whenever an independent structural theorem first makes $T$
similar to a contraction.

The hypothesis still sees only ordinary eigenvectors, not generalized
eigenvectors in long Jordan chains.  The 2024 construction hides its failure
of complete polynomial boundedness in large finite blocks.  Within each
block it places increasingly many distinct eigenvalues in a compact
pseudohyperbolic region, while the corresponding adjoint eigenlines are
orthogonal.  A bounded map cannot send almost parallel normalized Hardy
kernels to uniformly nonzero orthogonal vectors.  This simple observation
pinpoints why the strongest known candidate construction misses the source
hypothesis and why its long-chain obstruction cannot be imported verbatim.

\section{Two affirmative regimes}

\begin{theorem}[Completely polynomially bounded case]\label{thm:cpb}
Under the hypotheses of the source question, if $T$ is completely
polynomially bounded, then $T$ is similar to $S$.
\end{theorem}

\begin{proof}
Paulsen's theorem~\cite{Paulsen} gives an invertible $L$ for which
$R=LTL^{-1}$ is a contraction.  Set $X_0=XL^{-1}$.  Then
\[
 X_0R=SX_0,
 \qquad X_0^*u_\lambda=L^{-*}X^*u_\lambda.
\]
The map $X_0$ is a quasiaffinity, and
\[
 \frac{c}{\|L\|}\leq\|X_0^*u_\lambda\|
 \leq \|L^{-1}\|C \qquad(\lambda\in\D).
\]
Uchiyama's contraction theorem~\cite{Uchiyama} therefore gives $R\approx S$.
Since $T\approx R$, the conclusion follows.
\end{proof}

This identifies a sharp necessary feature of any negative example: it must
be polynomially bounded but not completely polynomially bounded.

\begin{theorem}[Unilateral weighted shifts]\label{thm:weighted}
Let $T e_n=w_ne_{n+1}$ be a unilateral weighted shift on $\ell^2(\mathbb
N_0)$ with every $w_n\neq0$.  If $T$ is power bounded and there is a
nonzero bounded $X$ satisfying $XT=SX$, then $T\approx S$.
In particular the source question has an affirmative answer for unilateral
weighted shifts, without using polynomial boundedness or
\eqref{eq:kernel-estimate}.
\end{theorem}

\begin{proof}
Put $P_0=1$ and $P_n=w_0\cdots w_{n-1}$ for $n\geq1$.  If $f=Xe_0$, then
intertwining gives
\[
 Xe_n=P_n^{-1}S^nf.
\]
Because $X$ is nonzero, $f\neq0$: if $Xe_0=0$, the displayed relation gives
$X=0$.  Consequently
\[
 |P_n|\geq\frac{\|f\|}{\|X\|}>0.
\]
Power boundedness supplies the reverse uniform estimate
$|P_n|=\|T^ne_0\|\leq\sup_k\|T^k\|$.  Hence the diagonal operator
$Le_n=P_ne_n$ is bounded and invertible.  The identity $TL=LS$ proves
$T\approx S$.
\end{proof}

\section{Interpolating-sequence rigidity}

Call $\Lambda\subset\D$ an $H^\infty$-interpolating sequence if every
bounded scalar family on $\Lambda$ is the restriction of some
$f\in H^\infty$.  Let $K_\Lambda$ denote an interpolation constant.

\begin{theorem}[Local Riesz rigidity]\label{thm:riesz}
Assume $T$ has polynomial bound $M$, $XT=SX$, and
\eqref{eq:kernel-estimate}.  For every $H^\infty$-interpolating sequence
$\Lambda$, the family
\[
 v_\lambda=\frac{X^*u_\lambda}{\|X^*u_\lambda\|},
 \qquad \lambda\in\Lambda,
\]
is a Riesz sequence.  Its Riesz constants depend only on $M$ and
$K_\Lambda$.  Moreover $X^*$ is bounded below on
$\clos\operatorname{span}\{u_\lambda:\lambda\in\Lambda\}$.
\end{theorem}

\begin{proof}
Fix a finite $F\subset\Lambda$ and signs
$\varepsilon_\lambda\in\{-1,1\}$.  For every $\eta>0$ there is a polynomial
$p$ with
\begin{equation}\label{eq:poly-interp}
 p(\lambda)=\varepsilon_\lambda\quad(\lambda\in F),
 \qquad \|p\|_\infty\leq K_\Lambda+\eta.
\end{equation}
Indeed, first interpolate by $f\in H^\infty$ with
$\|f\|_\infty\leq K_\Lambda$.  Radial dilation puts $f_r$ in the disk
algebra.  For $r$ sufficiently close to one, a finite Lagrange correction
makes its values on $F$ exact with norm at most $K_\Lambda+\eta/2$.
Uniform polynomial approximation followed by a second, arbitrarily small
Lagrange correction proves \eqref{eq:poly-interp}.

Since
\[
 T^*X^*u_\lambda=\overline\lambda X^*u_\lambda,
\]
$p(T)^*$ acts on $v_\lambda$, $\lambda\in F$, as the prescribed sign
change.  Its norm is at most $M(K_\Lambda+\eta)$.  Letting $\eta\downarrow0$
shows that every sign-change operator on the finite family has norm at most
$K_0=MK_\Lambda$.  It is its own inverse.  For scalars $a_\lambda$, averaging
over independent Rademacher signs gives
\[
 K_0^{-2}\sum_{\lambda\in F}|a_\lambda|^2
 \leq\left\|\sum_{\lambda\in F}a_\lambda v_\lambda\right\|^2
 \leq K_0^2\sum_{\lambda\in F}|a_\lambda|^2.
\]
Thus the $v_\lambda$ form a Riesz sequence.

The normalized Hardy kernels on an interpolating sequence are also a Riesz
sequence~\cite{Nikolski}.  Combining their upper Riesz estimate, the lower
estimate just proved, and $\|X^*u_\lambda\|\geq c$ shows that, on every finite
linear combination,
\[
 \left\|X^*\sum a_\lambda u_\lambda\right\|
 \geq c_\Lambda\left\|\sum a_\lambda u_\lambda\right\|
\]
for some $c_\Lambda>0$.  Taking closures proves the last assertion.
\end{proof}

The theorem does not by itself prove that $X^*$ is bounded below on all of
$\Htwo$.  A discrete interior interpolating sequence does not furnish a
sampling Riesz basis for $\Htwo$, and clustered kernels require information
about divided differences or generalized eigenvectors that is absent from
\eqref{eq:kernel-estimate}.

\section{The cluster--orthogonality obstruction}

Let
\[
 \rhoh(\lambda,\mu)=\left|\frac{\lambda-\mu}
 {1-\overline\mu\lambda}\right|
\]
be the pseudohyperbolic distance.

\begin{lemma}[Cluster--orthogonality]\label{lem:cluster}
Let $A:\Htwo\to\K$ be bounded.  Suppose that the vectors
$Au_\lambda$, $\lambda\in\Lambda$, lie in mutually orthogonal
one-dimensional subspaces and
\[
 \|Au_\lambda\|\geq c>0\qquad(\lambda\in\Lambda).
\]
Then $\Lambda$ is uniformly pseudohyperbolically separated.  Quantitatively,
for distinct $\lambda,\mu\in\Lambda$,
\begin{equation}\label{eq:separation}
 1-\sqrt{1-\rhoh(\lambda,\mu)^2}
 \geq \frac{c^2}{\|A\|^2}.
\end{equation}
\end{lemma}

\begin{proof}
The normalized Hardy kernels satisfy
\[
 |\langle u_\lambda,u_\mu\rangle|
 =\sqrt{1-\rhoh(\lambda,\mu)^2}.
\]
Choose $\omega\in\mathbb T$ so that
$\langle u_\lambda,\omega u_\mu\rangle$ is nonnegative.  Orthogonality of
the image lines gives
\[
 2c^2\leq\|A(u_\lambda-\omega u_\mu)\|^2
 \leq\|A\|^2\|u_\lambda-\omega u_\mu\|^2
 =2\|A\|^2\bigl(1-\sqrt{1-\rhoh(\lambda,\mu)^2}\bigr).
\]
This is \eqref{eq:separation}.
\end{proof}

\section{Audit of the 2024 candidate construction}\label{sec:audit}

Gamal's 2024 Theorem~6.3~\cite{Gamal2024} constructs finite-dimensional
blocks
\[
 T_N=\begin{pmatrix}T_{N1}&\Gamma_N\\0&T_{N0}\end{pmatrix}
 \approx T_{B_N}
\]
such that $\sup_N M_{pb}(T_N)<\infty$ but
$\sup_N M_{cpb}(T_N)=\infty$.  The lower diagonal contains
$d_N=2^{N+1}$ pairwise distinct eigenvalues
$\lambda_{N,l}$ in a fixed compact subdisk.  A common disk automorphism is
later applied to each block, and Theorem~7.1 combines the blocks into a
cyclic polynomially bounded operator quasisimilar to $T_B$ and not similar
to a contraction.  Corollary~2.3 then gives a polynomially bounded
quasiaffine transform of $S$ which is not similar to a contraction.

\begin{theorem}[The direct 2024 extension misses the kernel estimate]
\label{thm:audit}
For the Theorem~6.3 block sequence used in the Theorem~7.1/Corollary~2.3
shift extension of~\cite{Gamal2024}, there is no quasiaffinity
$Z$ intertwining the resulting polynomially bounded operator $\mathbf T$
to $S$ for which
\[
 \inf_{\lambda\in\D}\|Z^*u_\lambda\|>0.
\]
Thus this construction does not answer the source question negatively.
\end{theorem}

\begin{proof}
Write
\[
 T_{N0}=\bigoplus_{n=0}^N D_N+\mathbf S_N,
\]
where $D_N$ is diagonal in an orthonormal basis
$e_{N,1},\ldots,e_{N,d_N}$.  The chain subspaces generated by distinct
$e_{N,l}$ are mutually orthogonal.  Since the upper diagonal spectrum of
$T_N$ consists of distinct points $\nu_{N,j}$, none equal to any
$\lambda_{N,l}$, the triangular form gives
\[
 \ker(T_N^*-\overline{\lambda_{N,l}}I)
 =\{0\}\oplus\ker(T_{N0}^*-\overline{\lambda_{N,l}}I).
\]
Consequently these adjoint eigenlines are mutually orthogonal as $l$ varies.
The common disk automorphism used on block $N$ preserves both these lines
and all pseudohyperbolic distances.

Let $\zeta_{N,l}$ be the transformed eigenvalues.  The construction may,
and in the proof does, take the perturbation radii in a fixed compact
subdisk (one may decrease each radius without affecting any estimate).
Since $d_N\to\infty$, compactness of pseudohyperbolic balls implies
\begin{equation}\label{eq:cluster-pairs}
 \min_{l\neq k}\rhoh(\zeta_{N,l},\zeta_{N,k})\longrightarrow0
\end{equation}
along a subsequence.

The final $\mathbf T$ is polynomially bounded and $\mathbf T\prec S$.
For every $\zeta\in\D$, the standard defect-one theorem for polynomially
bounded quasiaffine transforms of $S$ gives
\[
 \dim\ker(\mathbf T^*-\overline\zeta I)=1
\]
(apply the result at zero after a disk automorphism; see
\cite{BercoviciPrunaru,Gamal2019}).  At a zero $\zeta_{N,l}$ of $B$, the
canonical Corollary~2.3 eigenvector has zero component in $BH^2$ and lies in
the finite-block eigenline above.  It is nonzero and hence spans the whole
eigenspace.  Therefore, for \emph{any} intertwiner $Z\mathbf T=SZ$,
$Z^*u_{\zeta_{N,l}}$ lies in that line.  The image lines are mutually
orthogonal for fixed $N$.

If the displayed uniform lower bound held, Lemma~\ref{lem:cluster} applied
to $A=Z^*$ would give one positive separation constant for every set
$\{\zeta_{N,l}:1\leq l\leq d_N\}$.  This contradicts
\eqref{eq:cluster-pairs}.
\end{proof}

For the canonical Corollary~2.3 intertwiner $Y$, the same obstruction is
visible in the exact identity
\[
 (1-|\lambda|^2)\|Y^*h_\lambda\|^2
 =|B(\lambda)|^2
 +(1-|\lambda|^2)\|Y_0^*h_{B,\lambda}\|^2.
\]
At every zero of $B$, the shift summand disappears completely and the norm
is forced onto the clustered finite-block eigenline.  The outer product
identity $X_0Y_0=g(T_B)$ from Theorem~7.1 does not repair this: its direct
lower estimate contains $|g(\lambda)|$, and the chosen outer $g$ tends
rapidly to zero at the shifted blocks.  A corona bound for $(B,g)$ would
instead make $g(T_B)$ invertible and force similarity of the base operator
to the model.

\section{Upgrade audit, limitations, and novelty status}

Eight focused routes were pursued.  Direct use of the 2024 examples led to
the kernel identity above.  The outer-product route failed at the necessary
corona bound.  Finite-block renormalization was ruled out by
Lemma~\ref{lem:cluster}.  Adding a coprime ``good'' model factor cannot help
at an actual zero of $B$, because that eigenvector belongs to the $B$ factor;
sharing the zero raises geometric multiplicity and is incompatible with a
quasiaffinity to a multiplicity-one shift.  Merging each finite block into
one very long cyclic Jordan chain preserves scalar polynomial estimates
under sufficiently small finite-dimensional perturbations, but the
intertwiner expands across the small connecting links and no uniform
eigenvector normalization was obtained.  The complete-polynomial-bound and
weighted-shift routes give Theorems~\ref{thm:cpb} and~\ref{thm:weighted}.
Finally, bounded analytic interpolation gives Theorem~\ref{thm:riesz}, but
there is no global interior-kernel sampling Riesz basis on which to close
the argument.

A bounded novelty search used the exact source title and the phrases
``polynomially bounded'', ``holomorphic eigenvectors'', and ``similar to the
unilateral shift'', together with later papers of the source author.  The
official arXiv records and TeX sources of arXiv:1901.03883 and
arXiv:2412.14130 were inspected.  The latter explicitly answers a different
question of K\'erchy but neither cites nor claims to resolve the 2019
eigenvector-norm question.  No later explicit resolution of the exact
question was found in this search.

Theorem~\ref{thm:audit} is an agent-derived implication from the published
block construction, not a claim made in~\cite{Gamal2024}.  The principal
human-review points are (i) the defect-one passage after disk automorphisms,
(ii) the orthogonality of the lower diagonal eigenlines in Theorem~6.3, and
(iii) the compact pseudohyperbolic packing step.  The packet claims only the
partial results above; the full polynomially bounded question remains open.

\begin{thebibliography}{9}

\bibitem{BercoviciPrunaru}
H.~Bercovici and B.~Prunaru,
\emph{Quasiaffine transforms of polynomially bounded operators},
Arch. Math. (Basel) 71 (1998), 384--387.

\bibitem{Gamal2019}
M.~F.~Gamal',
\emph{On power bounded operators with holomorphic eigenvectors, II},
Acta Sci. Math. (Szeged) 86 (2020), 549--562;
arXiv:1901.03883.

\bibitem{Gamal2024}
M.~F.~Gamal',
\emph{Examples of cyclic polynomially bounded operators that are not
similar to contractions}, arXiv:2412.14130 (2024).

\bibitem{Nikolski}
N.~K.~Nikolski,
\emph{Operators, Functions, and Systems: An Easy Reading}, Vols.~I--II,
Mathematical Surveys and Monographs 92, AMS, 2002.

\bibitem{Paulsen}
V.~I.~Paulsen,
\emph{Every completely polynomially bounded operator is similar to a
contraction}, J. Funct. Anal. 55 (1984), 1--17.

\bibitem{Uchiyama}
M.~Uchiyama,
\emph{Curvatures and similarity of operators with holomorphic eigenvectors},
Trans. Amer. Math. Soc. 319 (1990), 405--415.

\end{thebibliography}

\end{document}
